\section{Pro-objects, finite extensions, and transfer countability}
\label{sec:pro-transfer}

We now pass from the finite stable transfer image to the cohomology of an
individual finite marking level.  All coefficient fields in this section
are finite.  In particular, a countable-dimensional coefficient space is
a countable set, and a finite-dimensional coefficient space is a finite
set.  Both facts will be used.

Fix $1\leq m<n$ and retain the notation
\[
 A=A_m,\qquad B=A[1/x],\qquad B_U=eA'[1/x],\qquad
 \mathcal P=x^{-D_0}A'
\]
of \cref{sec:finite-markings,sec:transfers,sec:cutoff}.
Choose a cofinal decreasing sequence $(U_\nu)_{\nu\geq0}$ of normal
product congruence subgroups of the marking group $J$.  Discard finitely
many terms so that these groups are uniform, all the finite-isotropy
conclusions hold, and the transfer $S$ is defined.  The lattice
$\mathcal P$ and its whole finite free model may depend on $\nu$.
Write
\[
 X^i_\nu=H^i_{\mathrm{cts}}(G,B_{U_\nu}),\qquad
 M^i_\nu=H^i_{\mathrm{cts}}(G,\mathcal P_\nu).
\]
The $M^i_\nu$ are compact.  Throughout the section, the hypothesis on
the next height is precisely the coefficient-finiteness hypothesis in
\cref{lem:mark-strips}.  Thus every fixed whole-lattice shift has
countable localization kernel and cokernel.  No finiteness assertion at
the current height is assumed.

\begin{theorem}[Countability at fixed marking levels]
\label{thm:pro-countability}
Assume the next-height coefficient-finiteness hypothesis of
\cref{lem:mark-strips}.  There is an index $\nu_0$ such that
\[
 H^i_{\mathrm{cts}}(G,B_{U_\nu})
 \quad\text{is countable-dimensional over $\kk$}
\]
for every $\nu\geq\nu_0$ and every integer $i$.
\end{theorem}

We will use the following consequences of the preceding sections.
For every fixed $\nu$ and $N\geq1$, the map
\[
 S^NM^i_\nu\longrightarrow S^NX^i_\nu
\]
has countable kernel and cokernel.  To see this, first use
\cref{lem:mark-strips} for the whole model $A'[1/x]$.  The extension of
$S$ to that whole model begins with projection by $e$, so its positive
powers discard the complementary component.  Taking the images of the
two commuting $S^N$ maps preserves the assertion about countable kernel
and cokernel.  Furthermore, \cref{prop:cutoff-compact} says that if
\begin{equation}
\label{eq:pro-consecutive}
 S^NX^i_\nu/S^{N+1}X^i_\nu\quad\text{is countable-dimensional},
\end{equation}
then $S$ on $S^NX^i_\nu$ has countable kernel and cokernel, and
\begin{equation}
\label{eq:pro-finite-intersection}
 F^i_\nu:={\bigcap}_{a\geq0}
 \operatorname{im}\bigl(S^aM^i_\nu\longrightarrow X^i_\nu\bigr)
 \quad\text{is finite}.
\end{equation}
These assertions concern actual cohomology groups, without taking
closures of images.  The closedness needed in
\cref{eq:pro-finite-intersection} is part of
\cref{prop:cutoff-compact}.

For $m>1$, the proof has two steps. First we establish
\cref{eq:pro-consecutive} in each degree with an exponent independent
of the sufficiently deep marking level. The compact comparison above then
makes $F^i_\nu$ finite. This does not yet control all of $X^i_\nu$:
a finite stable intersection alone places no bound on the part lost
under iteration.

The second step rules out an uncountable remainder. The uniform
finite-isotropy bound leaves only finitely many cohomological degrees.
If countability fails, choose the largest degree $b$ in which
uncountable groups persist at arbitrarily deep marking levels. Put
$d=(m-1)(n-m)$, the dimension of the translation distribution algebra.
We will construct a finite $d$-fold translation extension which, for one fixed
root order $h_*$ and every sufficiently deep level $U_\nu$, induces an
operation
\[
 \varepsilon_\nu:
 H^{b-d}_{\cts}(G,B_{U_\nu}^{1/h_*})\longrightarrow X^b_\nu
\]
with countable cokernel. If $b-d<0$, its source is zero and countability
already follows.

Otherwise the next-height boundary hypothesis gives a subspace of the
source with countable quotient whose classes have positive-order cocycle
representatives. A second fixed $d$-fold extension, now a self-extension
of the trivial translation coefficient, supplies a boundary operator
$\mathcal B$ on ordinary marked cochains. The high-transfer comparison
will identify its coinduced operations, after inclusion from the fixed
root source and transfer back to $B_{U_\nu}$, with nonzero scalar
multiples of $\varepsilon_\nu$. This equality will hold for
every high transfer at each sufficiently deep fixed level.

For a positive representative $w$, powering drives its order to
infinity, while the boundary operator $\mathcal B$ loses only a fixed
pole order at that level. Its high-power boundary representatives
therefore lie in $\mathcal P_\nu$. The factorization then puts the
fixed class $\varepsilon_\nu[w]$ in every sufficiently late compact
transfer image, hence in $F^b_\nu$. This intersection is finite. Since
the positive-representative subspace has countable quotient, the whole
image of $\varepsilon_\nu$ is countable. Its countable cokernel now
contradicts the choice of $b$.

The category introduced next records this obstruction by making
countable spaces zero while retaining each sufficiently deep marking
level. Exact pro-completion will first control the transfer images and
then construct the top extension operation. A finite equivariant Ext
comparison will put all its high-transfer factorizations on one common
tail of actual marking levels. When $m=1$, the splitting
$S\Fr^{\log_p h_0}=1$ replaces the positive-dimensional translation
algebra; that case is treated at the end of the section.

\subsection{The category recording a sufficiently deep tail}

Let $\mathcal V$ be the category of $\kk$-vector spaces and let
$\mathcal V_{\leq\aleph_0}$ be its full subcategory of
countable-dimensional spaces.  This is a Serre subcategory: subspaces,
quotients, and extensions of countable-dimensional spaces are
countable-dimensional.  Put
\[
 \mathcal A_0=\mathcal V/\mathcal V_{\leq\aleph_0}.
\]
The quotient functor is exact.  A vector space becomes zero in
$\mathcal A_0$ exactly when it is countable-dimensional, and a linear
map becomes an isomorphism exactly when its kernel and cokernel are
countable-dimensional.

Define $\mathcal A$ to be the category of tail families in
$\mathcal A_0$.  Its objects are sequences $(Y_\nu)_{\nu\geq0}$ and
\[
 \operatorname{Hom}_{\mathcal A}(Y,Z)
 =\mathop{\operatorname{colim}}_{\nu_0}
   \prod_{\nu\geq\nu_0}
   \operatorname{Hom}_{\mathcal A_0}(Y_\nu,Z_\nu).
\]
Composition is componentwise after restricting to a common tail.
Kernels and cokernels are computed componentwise on a tail; a finite
diagram can always be placed on a common tail.  These descriptions
prove that $\mathcal A$ is abelian.  They also show that an object of
$\mathcal A$ is zero exactly when all its components vanish in
$\mathcal A_0$ after some index.  In particular, an equality in
$\mathcal A$ that involves a fixed finite diagram means an equality
modulo countable-dimensional spaces at all sufficiently deep actual
levels.  We do not replace these families by a direct limit of their
underlying vector spaces.

Suppose first that $m>1$, and set
\[
 \Lambda=\kk[[D_1,\ldots,D_d]],\qquad d=(m-1)(n-m)>0,
 \qquad \mathfrak m=(D_1,\ldots,D_d).
\]
Let $\mathcal C$ denote the abelian category of finite-length
$\Lambda$-modules.  For $E\in\mathcal C$, the translation torsor
defines the associated coefficient bundle $V_{E,U_\nu}$ for all
sufficiently large $\nu$.  Here a plain finite module can be given
trivial auxiliary action after shrinking $U_\nu$: its action factors
through a finite quotient of $\Lambda$, and a sufficiently deep
marking group acts trivially on that quotient.  The same choice works
for every arrow in a fixed finite diagram.  Consequently
\[
 \mathcal C^j(E)
   =\bigl(C^j_{\mathrm{cts}}(G,V_{E,U_\nu})\bigr)_\nu
 \quad\text{and}\quad
 T^i(E)=H^i\bigl(\mathcal C^\bullet(E)\bigr)
\]
are well-defined in $\mathcal A$, independently of omitted initial
terms.  By \cref{lem:cts-exact,prop:trans-roots}, each $\mathcal C^j$
is an exact additive functor.  Thus $T^\bullet$ has the long exact
sequences associated to short exact sequences in $\mathcal C$.
There is a uniform integer $c$ for which
\begin{equation}
\label{eq:pro-degree-bound}
 T^i(E)=0\qquad(i>c, E\in\mathcal C),
\end{equation}
by \cref{prop:mark-finite-isotropy,prop:cts-isotropy-bound}.
As usual, all these functors vanish in negative cohomological degrees.

Choose $q_0$ as in \cref{prop:trans-roots}, with all powers used below
fixing $\kk$, and write $q_a=q_0^a$.  The ideals
$(D_1^{q_a},\ldots,D_d^{q_a})$ are cofinal with the powers of
$\mathfrak m$.  Put
\[
 \Lambda_q=\Lambda/(D_1^q,\ldots,D_d^q).
\]
To identify the transition maps, use the perfect pairing of
\cref{sec:transfers} on the plain $\Lambda$-modules:
\[
 \alpha_q:\Lambda_q\xrightarrow{\sim}\mathsf N_q,
 \qquad
 \alpha_q(a)(b)=
 [D_1^{q-1}\cdots D_d^{q-1}](ab).
\]
Here the brackets extract the indicated monomial coefficient.
For allowed powers $q=st$, the regular quotient $\Lambda_q\to\Lambda_s$ becomes
under these pairings the coinduction of the translation augmentation
$\mathsf N_t\to\kk$ from the $s$-power subalgebra. Indeed, the
successive coefficient extractions select exponent
$(s-1)+s(t-1)=q-1$ in each variable. By
\cref{prop:trans-roots,prop:trans-cartier}, its coefficient-bundle map
is the corresponding root of a Cartier transfer. Thus the inverse
system of regular quotients records the classes that persist under
repeated transfer.

These pairings forget the auxiliary action; the natural actions on the
coinduced diagrams remain those of \cref{prop:trans-roots}. The
whole-diagram root identification, followed by powering on
coefficients, gives an isomorphism of inverse systems
\begin{equation}
\label{eq:pro-stationary}
 \{T^i(\Lambda_{q_a})\}_{a\geq0}
 \simeq
 \{T^i(\kk)\xleftarrow{S}T^i(\kk)
                    \xleftarrow{S}T^i(\kk)\xleftarrow{S}\cdots\}.
\end{equation}
The maps on the left are induced by the quotient maps of $\Lambda$.
Choose the scalar identifications recursively along this sequence so
that its transition maps are exactly $S$.  This is possible because
\cref{prop:trans-volume,prop:trans-cartier} make each discrepancy a
nonzero constant.  Each fixed finite part of this comparison is
realized on a sufficiently deep marking tail.  Thus
\cref{eq:pro-stationary} is an identity in
$\operatorname{Pro}(\mathcal A)$; it does not assert that the natural
actions on every $\mathsf N_q$ are simultaneously trivial at one
finite marking level.

\subsection{Exact completion and stabilization of the actual images}

We record the categorical facts with the precision needed for a
uniform stabilization index.

\begin{lemma}
\label[lemma]{lem:pro-constants}
For an abelian category $\mathcal B$, the constant objects in
$\operatorname{Pro}(\mathcal B)$ are closed under kernels and
cokernels of maps between them, and under extensions.
If a decreasing system of subobjects $J_a\subset X$ represents a
constant subobject $I\subset X$ in
$\operatorname{Pro}(\mathcal B)$, then $J_a=I$ for all sufficiently
large $a$.
\end{lemma}

\begin{proof}
The embedding by constant objects is fully faithful and exact, which
gives the kernel and cokernel statements.  For extensions it is enough
to give the dual argument in $\operatorname{Ind}(\mathcal B)$.
Suppose
\[
 0\longrightarrow A\longrightarrow E\longrightarrow B
 \longrightarrow0
\]
is exact with $A$ and $B$ constant.  Write $E$ as a filtered system
of objects $E_j$ of $\mathcal B$.  The images of the maps
$E_j\to B$ are objects of $\mathcal B$, and their filtered colimit
is $B$.  The defining Hom formula for ind-objects says that a map
from a constant object to a filtered colimit factors through one
term.  Applied to $\operatorname{id}_B$, it shows that
$E_j\to B$ is an epimorphism for some $j$.  Let
$K_j=\ker(E_j\to B)$.  The induced map $K_j\to A$ is a map in
$\mathcal B$.  The pushout of
$K_j\hookrightarrow E_j$ along $K_j\to A$ is $E$: the natural map
from that pushout has identity maps on the kernel $A$ and quotient
$B$, hence is an isomorphism.  This pushout is a cokernel between
objects of $\mathcal B$, so is constant.  Dualizing proves extension
closure in the pro-category.

For the second assertion, the morphism from the constant object $I$
to $\{J_a\}$ is a compatible family of maps $I\to J_a$.  Since its
composite with $J_a\hookrightarrow X$ is the fixed inclusion, these
are inclusions identifying $I$ as a subobject of every $J_a$.
The system $\{J_a/I\}$ is pro-zero.  Thus for any chosen $a$ some
later transition $J_b/I\to J_a/I$ is zero.  This transition is
monic, so $J_b/I=0$.  All still later subobjects contain $I$ and are
contained in $J_b=I$, proving the assertion.
\end{proof}

For a finitely generated $\Lambda$-module $L$, let
\[
 \widehat L_{\mathrm{pro}}
       =\{L/\mathfrak m^aL\}_{a\geq1}
       \in\operatorname{Pro}(\mathcal C).
\]
This completion functor is exact.  Indeed, for
$0\to L'\to L\to L''\to0$, the levelwise exact sequences use
$L'/(L'\cap\mathfrak m^aL)$ as their first terms.  Artin--Rees
makes the filtrations $L'\cap\mathfrak m^aL$ and
$\mathfrak m^aL'$ cofinal, identifying these first terms with
$\widehat L'_{\mathrm{pro}}$.
This is the same Artin--Rees and cofinal-filtration argument that
proves exactness of ordinary completion in \cite[Tag 00MA]{stacks};
here it identifies the completed sequences as pro-objects.

An exact functor between abelian categories extends exactly to their
pro-categories: after a common reindexing, a finite diagram has
levelwise kernels and cokernels, and the functor preserves these.
Apply this to each $\mathcal C^j$.  Define
\[
 \overline{\mathcal C}^{\,\bullet}(L)
      =\mathcal C^\bullet(\widehat L_{\mathrm{pro}}),
 \qquad
 \overline T^{\,i}(L)
      =H^i\bigl(\overline{\mathcal C}^{\,\bullet}(L)\bigr).
\]
Thus
$\overline T^{\,i}(L)=\{T^i(L/\mathfrak m^aL)\}_a$ in
$\operatorname{Pro}(\mathcal A)$, and short exact sequences of
finitely generated modules give long exact sequences.  If $L$ has
finite length, its completion is constant, so
$\overline T^{\,i}(L)$ is the constant object $T^i(L)$.

\begin{proposition}[Uniform image stabilization]
\label[proposition]{prop:pro-stabilization}
For every $i$, the pro-object $\overline T^{\,i}(\Lambda)$ is
constant, and its canonical projection to $T^i(\kk)$ is monic.
It is annihilated by $\mathfrak m$.  Moreover, there is an integer
$N_i\geq1$, independent of the sufficiently deep marking level,
such that
\[
 S^{N_i}X^i_\nu/S^{N_i+1}X^i_\nu
 \quad\text{is countable-dimensional for all sufficiently large $\nu$}.
\]
\end{proposition}

\begin{proof}
We descend on $i$, starting above the bound in
\cref{eq:pro-degree-bound}.  Assume the assertion of constancy has
been proved for $\overline T^{\,j}(\Lambda)$ with $j>i$.
The ideal $\mathfrak m$ has a finite resolution by finite free
$\Lambda$-modules, given by the positive part of the Koszul
resolution of $\kk$.  Induction on resolution length shows that
$\overline T^{\,j}(L)$ is constant for every $j>i$ and every $L$
with a finite free resolution.  The case of a free $L$ is the
descending induction hypothesis and finite additivity.  For
$0\to L'\to F\to L\to0$ with $F$ finite free and $L'$ of shorter
resolution length, the long exact sequence gives
\begin{multline*}
0\longrightarrow
\coker\bigl(\overline T^{\,j}(L')\to\overline T^{\,j}(F)\bigr)
\longrightarrow\overline T^{\,j}(L)\\
\longrightarrow
\ker\bigl(\overline T^{\,j+1}(L')\to\overline T^{\,j+1}(F)\bigr)
\longrightarrow0.
\end{multline*}
For $j>i$ the outer terms are constant by the induction on length
and \cref{lem:pro-constants}; extension closure gives the middle
term.  Taking $L=\mathfrak m$ and $j=i+1$ proves the required
constancy of $\overline T^{\,i+1}(\mathfrak m)$.

The exact sequence $0\to\mathfrak m\to\Lambda\to\kk\to0$
now identifies the image of
\[
 \overline T^{\,i}(\Lambda)\longrightarrow T^i(\kk)
\]
with the kernel of a map between the two constant objects
$T^i(\kk)$ and $\overline T^{\,i+1}(\mathfrak m)$.
Call this constant subobject $I^i$.  Under
\cref{eq:pro-stationary}, the image system is
\[
 J_a=\operatorname{im}\bigl(S^a:T^i(\kk)\to T^i(\kk)\bigr),
 \qquad J_{a+1}\subset J_a.
\]
The second part of \cref{lem:pro-constants} supplies one finite
$N_i$ with $J_a=I^i$ for all $a\geq N_i$.  Increase $N_i$ to at
least one.  Equality of these subobjects in $\mathcal A$ says
exactly that \cref{eq:pro-consecutive} holds for this $N_i$ at
every sufficiently deep actual level.  Its depth is uniform
because only this one consecutive-image comparison is involved.

By \cref{prop:cutoff-compact}, the endomorphism induced by $S$ on
$J_{N_i}=I^i$ has countable kernel and cokernel at these levels.
It is therefore an automorphism in $\mathcal A$.  Replacing the
stationary inverse system by its $N_i$th images gives an isomorphic
pro-object.  This replacement has underlying object $I^i$ and
invertible transition $S|_{I^i}$, so is constant.  Explicitly, its
map from the constant object $I^i$ to level $a$ is
$(S|_{I^i})^{-a}$ followed by inclusion.  At the initial level the
projection is the inclusion $I^i\hookrightarrow T^i(\kk)$.
This proves the claimed monomorphism and completes the descending
induction.

Finally, multiplication by any element of $\mathfrak m$ followed
by $\Lambda\to\kk$ is zero.  The resulting endomorphism of
$\overline T^{\,i}(\Lambda)$ has zero composite with the monic
projection just proved.  It is zero.  Thus $\mathfrak m$ annihilates
the pro-object.
\end{proof}

\subsection{Finite Yoneda extensions and the top Koszul operation}

We spell out why operations computed with completed free modules
can be realized by finite coefficient diagrams.

\begin{lemma}[Finite-length extension comparison]
\label[lemma]{lem:pro-finite-ext}
For finite-length $\Lambda$-modules $E,F$, the natural map
\[
 \operatorname{Ext}^j_{\mathcal C}(E,F)
       \longrightarrow \operatorname{Ext}^j_\Lambda(E,F)
\]
is an isomorphism for every $j$.  Here the left side is Yoneda Ext:
its classes are represented by exact diagrams whose middle terms
all have finite length.
\end{lemma}

\begin{proof}
Let $\mathcal D$ be the category of $\mathfrak m$-power-torsion
$\Lambda$-modules.  Every finitely generated submodule of an object
of $\mathcal D$ has finite length: finitely many generators are
killed by one power of $\mathfrak m$, and the corresponding
quotient ring is Artinian.  Thus
$\mathcal D=\operatorname{Ind}(\mathcal C)$.

Finite endpoints have the same Yoneda extensions in these two
categories.  One can see the finite realization directly.  Starting
at the right endpoint of an extension in $\mathcal D$, choose
finitely many lifts of its generators, then take the finite-length
submodule that they generate.  Its kernel is finite-length.  Lift
generators of that kernel into the preceding middle term and
continue.  At the last step include the whole left endpoint.
This constructs a finite-length subextension representing the
given class.  The same construction can be made with a prescribed
finite subcomplex: include its terms before choosing the next
lifts.  To check injectivity, take a finite Yoneda zigzag witnessing
equality.  Its directed underlying path has no directed cycles.
Process its vertices in an order compatible with the arrows,
choosing a finite subextension that contains the images of every
already chosen predecessor; at a finite original vertex retain the
whole extension.  All arrows then restrict to a finite zigzag.
Hence the map on Ext is both surjective and injective.
Equivalently, this is extension fullness of an abelian
category in its ind-completion
\cite[Theorem 2.2.1]{coulembier2020}.

For comparison with all $\Lambda$-modules, use that the
$\mathfrak m$-power-torsion submodule of an injective module is
itself injective over the Noetherian ring $\Lambda$
\cite[Lemma 47.3.9(2), Tag 08XW]{stacks}.
If $M\in\mathcal D$ embeds in an ordinary injective $I$, that
embedding factors through $I[\mathfrak m^\infty]$.  Its cokernel
is still $\mathfrak m$-power torsion.  Iterating this construction
gives an ordinary injective resolution entirely in $\mathcal D$.
Each term is also injective in the full abelian subcategory
$\mathcal D$, so the resolution computes both Ext groups, proving
$\operatorname{Ext}_{\mathcal D}^j(E,F)
=\operatorname{Ext}_\Lambda^j(E,F)$ and the lemma.
\end{proof}

An exact finite Yoneda diagram determines an operation on $T^\bullet$
by the composite of its connecting maps.  Functoriality of long exact
sequences and invariance under comparisons make this operation depend
only on its Ext class.

\begin{proposition}[A surjective top operation]
\label[proposition]{prop:pro-top-operation}
Suppose $T^\bullet(\kk)$ is not zero, and let $b$ be its largest
nonzero degree.  There is an allowed power index $s=q_0^{a_s}$ and a generator
\[
 e_s\in\operatorname{Ext}^d_\Lambda(\mathsf N_s,\kk)
\]
whose operation is an epimorphism in $\mathcal A$:
\[
 e_s:T^{b-d}(\mathsf N_s)\longrightarrow T^b(\kk).
\]
\end{proposition}

\begin{proof}
A composition series shows that $T^j(E)=0$ for every finite-length
$E$ and $j>b$.  In particular
$\overline T^{\,b+1}(\mathfrak m)=0$.  The exact sequence for
$0\to\mathfrak m\to\Lambda\to\kk\to0$, together with
\cref{prop:pro-stabilization}, gives an isomorphism
\begin{equation}
\label{eq:pro-top-identification}
 Y:=\overline T^{\,b}(\Lambda)\xrightarrow{\ \sim\ }T^b(\kk).
\end{equation}

For an allowed index $s$, identify $\mathsf N_s$ with
$\Lambda_s$ as a plain $\Lambda$-module.  Let $K_s^\bullet$
be the Koszul resolution on $D_1^s,\ldots,D_d^s$, in degrees
$[-d,0]$.  Apply exact completion and then the exact cochain-term
functors $\mathcal C^j$ to this bounded complex.  The augmented
horizontal rows remain exact, so the total complex computes the
constant object $\mathcal C^\bullet(\Lambda_s)$.
Its vertical-cohomology spectral sequence has
\[
 E_1^{a,j}(s)=
 \overline T^{\,j}(\Lambda)^{\binom d{-a}},
 \qquad -d\leq a\leq0,
 \qquad
 E_1^{a,j}(s)\Longrightarrow T^{a+j}(\Lambda_s).
\]
The finite column interval and \cref{eq:pro-degree-bound} give a
finite filtration in every total degree.  By
\cref{prop:pro-stabilization}, $\mathfrak m$ annihilates every
entry on this first page, so all first differentials vanish.

We compare these operations in the forward direction. Under the
pairings $\alpha_s$ above, the natural inclusion
$\mathsf N_s\to\mathsf N_{s'}$, dual to the regular quotient, is
\[
 \Lambda_s\longrightarrow\Lambda_{s'},\qquad
 1\longmapsto\prod_{j=1}^dD_j^{s'-s}
 \qquad(s'>s).
\]
Its Koszul comparison preserves the top exterior term while making
every other column map vanish on vertical cohomology. To see this,
lift the displayed injection to a map
$K_s^\bullet\to K_{s'}^\bullet$.
On the exterior-basis summand indexed by
$I\subset\{1,\ldots,d\}$, that map is multiplication by
\[
 \prod_{j\notin I}D_j^{s'-s}.
\]
The Koszul differential removes one index, and this formula
commutes with it because the missing factor changes $D_j^s$
to $D_j^{s'}$.  On the first spectral-sequence page the comparison
is the identity in column $-d$ and zero in every other column.
The latter maps remain zero on every subsequent page.

There is no incoming differential at $(-d,b)$.  Starting with
$E_2^{-d,b}=Y$, increase $s$ once for each possible page
$r=2,\ldots,d$.  Inductively suppose the whole object $Y$
survives to page $r$ at the current index.  At a larger index the
comparison remains the identity on this whole source.  Naturality
of $d_r$ and the zero comparison on its target column show that
the outgoing $d_r$ at the larger index is zero on all of $Y$.
The earlier differentials are still zero for the same reason.
After these finitely many increases,
$E_\infty^{-d,b}=Y$.  The leftmost-column edge of the finite
filtration is therefore an epimorphism
\[
 T^{b-d}(\Lambda_s)\twoheadrightarrow Y.
\]
When $d=1$ there are no pages to eliminate, and the same conclusion
holds directly.

The edge followed by \cref{eq:pro-top-identification} is the
operation of the chain map
\[
 K_s^\bullet\longrightarrow\Lambda[d]
                       \longrightarrow\kk[d]
\]
which projects the top exterior term onto $\Lambda$ and then
augments.  Applying $\operatorname{Hom}_\Lambda(-,\kk)$ to
$K_s^\bullet$ gives zero differentials; the displayed map is
therefore a generator of
$\operatorname{Ext}^d_\Lambda(\Lambda_s,\kk)\simeq\kk$.
By \cref{lem:pro-finite-ext} it has a finite-length Yoneda
representative.  To verify that this representative gives the
same operation, compare the two representing morphisms in the
bounded derived category of finitely generated $\Lambda$-modules,
using their finite free resolutions.  They agree there because
they represent the same Ext class.  Exact completion and the
exact cochain-term functors preserve quasi-isomorphisms of these
bounded complexes, as is seen by taking the exact horizontal
rows of their cones.  Their induced operations thus agree after
totalization.  This identifies the edge map with the operation
of the finite diagram and proves the proposition.
\end{proof}

\subsection{One auxiliary shrink for every high transfer}

We next establish the uniformity that is not supplied merely by
forming a germ of marking levels.  Let $T$ be the translation group
scheme, so that finite rational $T$-representations are the objects
of $\mathcal C$.  For a compact open subgroup $U\subset J$, let
$\mathcal R_U$ be the category of representations of $T\rtimes U$
which are rational for $T$ and discrete continuous for $U$.
The compatibility is the usual semidirect-product compatibility.

\begin{lemma}[Finite forgetting kernel]
\label[lemma]{lem:pro-common-shrink}
Let $U$ be a uniform pro-$p$ marking subgroup, and let $E,F$ be
finite objects of $\mathcal R_U$.  Each group
$\operatorname{Ext}^j_{\mathcal R_U}(E,F)$ is finite.  There is
one open subgroup $U'\subset U$ on which every element of the
kernel of
\[
 \operatorname{Ext}^j_{\mathcal R_U}(E,F)
      \longrightarrow\operatorname{Ext}^j_T(E,F)
\]
becomes zero.
\end{lemma}

\begin{proof}
There is an extension-descent spectral sequence
\begin{equation}
\label{eq:pro-equivariant-ext}
 H^a_{\mathrm{cts}}\bigl(U,
       \operatorname{Ext}^b_T(E,F)\bigr)
 \Longrightarrow\operatorname{Ext}^{a+b}_{\mathcal R_U}(E,F).
\end{equation}
Here is a construction verifying the coefficient-category
requirements.  Regard $U$ as the proconstant affine group scheme
whose coordinate algebra is the algebra of locally constant
$\kk$-valued functions on $U$.  The underlying vector space of
the coordinate coalgebra of $T\rtimes U$ is
$\mathcal O(T)\otimes C_{\mathrm{lc}}(U,\kk)$.
A cofree representation restricts to a cofree $T$-representation:
the semidirect action is changed to ordinary translation on the
$T$ coordinate by conjugating that coordinate by the retained
$U$ coordinate.  Its $T$-invariants are a cofree discrete
$U$-representation.  Such $U$-representations are acyclic for
invariants, since the locally constant coinduction functor is
exact; lifts through a vector-space surjection are chosen on the
finite image of a locally constant function.  Cofree resolutions
therefore give the spectral sequence for the composite of the
two invariant functors.  Tensoring with the dual of the finite
representation $E$ identifies that spectral sequence with
\cref{eq:pro-equivariant-ext}.

By \cref{lem:pro-finite-ext}, the translation Ext groups are the
ordinary $\Lambda$-Ext groups of finite-length modules.  The
Koszul resolution makes $\operatorname{Ext}^b_\Lambda(\kk,\kk)$
finite-dimensional and zero for $b>d$.  Composition series
and the long exact sequences in each variable give the same
two conclusions for all finite-length endpoints.
Their induced $U$-actions are discrete continuous.  Indeed a
finite Yoneda representative of a translation Ext class becomes
equivariant after shrinking $U$ to act trivially on its finite
endpoints and on the finite translation quotient through which
its middle terms act; give those middle terms trivial auxiliary
action.  This open subgroup fixes the class.  Finite-type
resolutions for the analytic group $U$ then make each group on the left of
\cref{eq:pro-equivariant-ext} finite-dimensional.  Only finitely
many terms contribute in a fixed total degree.  The abutment is
therefore finite-dimensional over the finite field $\kk$, hence
is a finite set.

Every representation of the indicated coalgebra is a union of
finite-dimensional subrepresentations.  As in the finite
subdiagram construction in \cref{lem:pro-finite-ext}, an
extension with finite endpoints, and any Yoneda comparison
between such extensions, can be realized by finite diagrams.
Fix a class in the displayed forgetting kernel.  Choose a
finite equivariant diagram representing it.  Its underlying
$T$-extension is zero, so there is a finite $T$-linear Yoneda
comparison with a representative of zero.  Shrink $U$ so that
it acts trivially on all the finite modules of the original
diagram.  Shrink further so that its action on the finite
translation quotients through which the comparison modules act
is trivial.  Give those comparison modules trivial auxiliary
action.  Every comparison arrow is then equivariant, and the
class has become zero.  Finally the forgetting kernel is a
finite set.  Intersect the finitely many open subgroups just
obtained for its elements.  This one intersection kills the
whole kernel.  A sufficiently deep normal product congruence
subgroup is contained in it.
\end{proof}

At a fixed marking level $U$, finite diagrams with their given
$T\rtimes U$-actions descend without trivializing those actions
at $U$.  Indeed, choose $\ell$ through which their translation
actions factor and a normal open $U'\subset U$ fixing both
$T/[p^\ell]T$ and every term.  The subgroup
$[p^\ell]T\rtimes U'$ is normal in $T\rtimes U$: the translation
subgroup is $U$-stable, $U'$ is normal in $U$, and $U'$ acts
trivially on $T/[p^\ell]T$.  The quotient of the full marked
lift torsor by $[p^\ell]T\rtimes U'$ is a torsor under the
finite group scheme
$\bigl(T/[p^\ell]T\bigr)\rtimes(U/U')$ over $\Spec B_U$,
by \cref{sec:finite-markings,prop:trans-roots}.
It is therefore a finite faithfully flat cover of that fixed base.
Here $U'$ specifies a cover over $B_U$.
Faithfully flat descent of the vector-space diagram tensored
with that cover gives its finite locally free associated bundles,
arrows, and exactness over $B_U$, with the commuting $G$-action.
This applies also to finite equivariant Yoneda comparisons.
Hence equality in equivariant Ext gives equality of the induced
operations on the actual cohomology of the associated bundles
at that level.  Restriction to any smaller $U$ preserves it.

\begin{lemma}[A fixed extension through all transfers]
\label[lemma]{lem:pro-high-transfer}
Fix the index $s$ and generator $e_s$ of
\cref{prop:pro-top-operation}, and choose a generator
$e'_0\in\operatorname{Ext}^d_\Lambda(\kk,\kk)$.
There is one sufficiently deep marking subgroup $U_*$ such
that for every allowed $q\geq s$ the composite
\begin{equation}
\label{eq:pro-composite-extension}
 \mathsf N_s\longrightarrow\mathsf N_q
 \xrightarrow{\operatorname{Coind}_q(e'_0)}
 \mathsf N_q[d]\longrightarrow\kk[d]
\end{equation}
equals a nonzero scalar multiple of $e_s$ in the equivariant
Ext group for $T\rtimes U_*$.  The first and last maps use the
natural auxiliary actions, which are retained for every $q$.
\end{lemma}

\begin{proof}
Choose finite Yoneda diagrams for $e_s$ and $e'_0$.
After one initial shrink $U_0$, their plain module diagrams
are equivariant with trivial auxiliary actions, agreeing with
the natural action at the fixed endpoint $\mathsf N_s$.
The power-subalgebra identification commutes with marking
changes by \cref{prop:trans-roots}, so coinduction gives
equivariant diagrams for every $q$ without making the action
on $\mathsf N_q$ trivial.  The natural inclusion
$\mathsf N_s\to\mathsf N_q$ is equivariant.  The space of
translation maps $\mathsf N_q\to\kk$ is one-dimensional;
the pro-$p$ group $U_0$ acts trivially on this line because
$\kk^\times$ has order prime to $p$.  Choose a nonzero map
on that line for the last arrow.

After forgetting $U_0$, the composite in
\cref{eq:pro-composite-extension} is a generator.  This can
be checked directly on the power-subalgebra free pairing.
With
$\Lambda^{(q)}=\kk[[D_1^q,\ldots,D_d^q]]$, the coefficient
of $\prod D_i^{q-1}$ gives the perfect pairing identifying
coinduction from $\Lambda^{(q)}$ with extension of scalars.
Extension of scalars carries the Koszul representative of
$e'_0$ to the top projection in the Koszul resolution on
$D_1^q,\ldots,D_d^q$.  Under the regular-module
identifications, the inclusion from the smaller block is
multiplication by $\prod D_i^{q-s}$ up to a unit.  The
Koszul comparison written in the proof of
\cref{prop:pro-top-operation} has top coefficient one.
Following with augmentation therefore leaves a nonzero top
coefficient in $\kk$.  Changes in the perfect-pairing
generators only multiply that coefficient by a unit, whose
augmentation is nonzero.  Thus the ordinary Ext class is
$c_qe_s$ for some $c_q\in\kk^\times$.

Subtract $c_qe_s$ from each composite, now in the fixed group
$\operatorname{Ext}^d_{\mathcal R_{U_0}}(\mathsf N_s,\kk)$.
All differences lie in its forgetting kernel.  Apply
\cref{lem:pro-common-shrink} to this fixed pair of endpoints.
One subgroup $U_*$ kills the entire kernel, and hence every
difference, independently of $q$.
\end{proof}

\subsection{The fixed boundary estimate and the contradiction}

We make explicit the cocycle estimate used after the common shrink.
Fix a marking level $U\subset U_*$ and one finite Yoneda diagram
for $e'_0$.  Applying its associated-bundle functor gives an exact
diagram of finite locally free modules over the affine ring $B_U$.
Each successive surjection in that diagram has a $B_U$-linear
section, since its quotient is finite projective.  Choose such
sections once.  Successively lifting a cocycle by these sections
and taking group differentials gives a boundary operator
\[
 \mathcal B:
 Z^j_{\mathrm{cts}}(G,B_U)
       \longrightarrow Z^{j+d}_{\mathrm{cts}}(G,B_U)
\]
representing $e'_0$ on cohomology.  It is an additive
$\kk$-linear operator; it need not be a morphism of
$B_U$-modules.

There is a constant $C_U$ such that $\mathcal B$ loses at most
$C_U$ of whole-lattice $x$-order, in every one of the bounded
degrees used below.  To verify the assertion, realize the finitely
many bundles as summands of finite free modules and choose
finitely many whole lattices in these realizations.  Every chosen
section is a finite matrix with entries having bounded poles.
The compact group action has a uniform bound on each of these
lattices.  Moreover $g(x)/x$ is an integral unit, so shifting a
lattice by $x^a$ changes none of these bounds.  A group
differential is a finite sum of action and face maps; hence
it loses at most a fixed additional pole order.  Adding the
losses of the finitely many sections and differentials gives
$C_U$.  Continuity of these operations on bounded-step
cochains follows from \cref{lem:cts-exact,lem:cts-comparison}.

For $q=q_0^a\geq s$, put $h=q^{m-1}$.
The entire-diagram identity in \cref{prop:trans-roots} transports
these fixed sections to the coinduced diagram.  Thus
\cref{lem:pro-high-transfer} yields, for a cocycle $w$ at the
fixed root level $h_s=s^{m-1}$,
\begin{equation}
\label{eq:pro-boundary-transfer}
 e_s[w]
   = c_{q,U}^{-1}
       S^a\bigl[\mathcal B(w^h)\bigr],
       \qquad c_{q,U}\in\kk^\times.
\end{equation}
This is equality in the actual group $H^{j+d}(G,B_U)$.
Here $w$ is first included from $B_U^{1/h_s}$ into
$B_U^{1/h}$, and powering identifies the latter root
coefficient diagram with the original one.  Since $h$ fixes
$\kk$ and commutes with $G$, $w^h$ is an ordinary cocycle.
The transfer after taking the boundary is $S^a$ under the
same powering identification.  In particular the bound on
the right of \cref{eq:pro-boundary-transfer} is the bound
for the one fixed operator $\mathcal B$, not a new bound
for each increasingly large coinduced diagram.

We also need positive representatives in the source.  For any
fixed integer $L>0$, \cref{lem:mark-strips} says that the image
of
\[
 H^j_{\mathrm{cts}}(G,x^LA')\longrightarrow
 H^j_{\mathrm{cts}}(G,B_U)
\]
has countable cokernel; the map includes the whole model and
then projects by $e$.  The idempotent has a fixed pole cost.
Powering by $h_s$ is a $\kk$-linear, $G$-equivariant
isomorphism from $B_U^{1/h_s}$ to $B_U$.  Choose $L$ beyond
the idempotent pole cost and pull the preceding image back
under this isomorphism.  We obtain a subspace
\begin{equation}
\label{eq:pro-positive-source}
 Z_U^j\subset H^j_{\mathrm{cts}}(G,B_U^{1/h_s})
\end{equation}
with countable quotient, whose classes have cocycle
representatives of strictly positive $x$-order on the
normalized root basis.  Increasing $L$ prescribes any fixed
positive lower bound on that order.  Whole-lattice
representatives are used before projection, so this assertion
does not specialize the idempotent across the boundary.

\begin{proof}[Proof of \cref{thm:pro-countability} for $m>1$]
If all $T^i(\kk)$ vanish, their finitely many potentially
nonzero degrees in \cref{eq:pro-degree-bound} have a common
vanishing tail, which is the desired conclusion.  Otherwise
choose the largest nonzero degree $b$ and the epimorphism
$e_s$ from \cref{prop:pro-top-operation}.
If $b-d<0$, its source is zero and already contradicts
$T^b(\kk)\ne0$.  We may therefore assume $j=b-d\geq0$.

Shrink the marking tail so that
\cref{lem:pro-high-transfer} holds and the consecutive-image
conclusion of \cref{prop:pro-stabilization} holds in degree
$b$.  Fix any one level $U=U_\nu$ in this tail.  The set
$F^b_\nu$ of \cref{eq:pro-finite-intersection} is finite.
Take a class in the positive subspace
\cref{eq:pro-positive-source} and a cocycle representative
$w$ as described there.  If its positive order is at least
$\epsilon>0$, its $h$th powers have whole-basis order tending
to positive infinity as $h=q^{m-1}$ increases.  Indeed first
power out the fixed root denominator $h_s$.  The resulting
whole finite model has integral multiplication laws, so
further Frobenius powers multiply the lower order bound;
the fixed projection and basis changes contribute only their
fixed pole costs.  The estimate for $\mathcal B$ therefore
puts $\mathcal B(w^h)$ in $\mathcal P_\nu$ for all
sufficiently large $a$.  Regard this cochain in the whole
localized model through $B_U=eB'\subset B'$.  Its differential
vanishes there because it vanishes in $B_U$.  Since
$\mathcal P_\nu$ is $G$-stable and
$\mathcal P_\nu\hookrightarrow B'$ is injective, it is a
$\mathcal P_\nu$-valued cocycle.

By \cref{eq:pro-boundary-transfer}, the fixed class $e_s[w]$
then belongs to
\[
 \operatorname{im}\bigl(S^aM^b_\nu\to X^b_\nu\bigr)
\]
for every sufficiently large $a$.  Multiplication by the
nonzero constant in that equation does not change this
subspace.  These image subspaces are decreasing, so membership
in every sufficiently late one implies membership in their
entire intersection $F^b_\nu$.  The threshold may depend on
$w$; no uniform threshold over all cocycles is needed.

Thus $e_s$ maps $Z_U^j$ into a finite-dimensional space.
The quotient of its source by $Z_U^j$ is countable-dimensional,
so the whole image of $e_s$ at this fixed level is
countable-dimensional.  This holds at every level in the
chosen tail.  Its morphism in $\mathcal A$ is therefore zero.
It was an epimorphism onto the nonzero object $T^b(\kk)$,
a contradiction.  Hence all $T^i(\kk)$ vanish, and the
uniform degree bound gives one common tail for all degrees.
\end{proof}

\subsection{The multiplicative-type case}

\begin{proof}[Proof of \cref{thm:pro-countability} for $m=1$]
By \cref{prop:trans-cartier}, the normalized transfer satisfies
\[
 S\Fr^{\log_p h_0}=1
\]
on the open coefficient algebra and on its group cochains.
It is therefore surjective on every $X^i_\nu$.  In
particular \cref{eq:pro-consecutive} holds, and
\cref{eq:pro-finite-intersection} gives a finite intersection
$F^i_\nu$ at each sufficiently deep marking level.

Choose a sufficiently positive whole-lattice shift $x^LA'$,
with $L$ larger than the fixed projection pole cost.  The
image of its cohomology in $X^i_\nu$ has countable cokernel,
by \cref{lem:mark-strips}.  If a class $z$ is represented
by a projected cocycle from this lattice, its iterated
Frobenius powers are integral in the whole model after
projection: their order grows, while the idempotent and basis
costs are fixed.  They consequently lie in $\mathcal P_\nu$
for every sufficiently large power.  The identity splitting
Frobenius gives
\[
 [z]=S^a\bigl[z^{h_0^a}\bigr]
\]
in the actual open cohomology.  As before, the image
subspaces are decreasing, so this class belongs to
$F^i_\nu$.  The image of the positive-lattice cohomology is
therefore finite; its cokernel is countable.  Hence
$X^i_\nu$ is countable.  The uniform cohomological bound
and the finitely many required choices give a common
sufficiently deep tail for all degrees.
\end{proof}
